\subsubsection{Model bazowy (Aligned\_Pretrained)}
Model bazowy trenowany jest jako klasyfikator z głowicą ArcMargin.

\begin{table}[h]
\centering
\small
\begin{tabular}{ll}
\toprule
\textbf{Hiperparametr} & \textbf{Wartość} \\
\midrule
Rozmiar partii & 32 \\
Epoki & 25 \\
Optymalizator & SGD (momentum=0.9) \\
LR & $1 \times 10^{-2}$ \\
Zanik wag & brak (nie ustawiono) \\
Wczesne zatrzymanie & Tak (cierpliwość=5) \\
Walidacja & 500 par z testu (średnia cos) \\
Augmentacja okluzji & $p=0.7$, wysokość 20 px \\
\bottomrule
\end{tabular}
\caption{Hiperparametry modelu bazowego (Aligned\_Pretrained).}
\label{tab:baseline_config}
\end{table}

\subsubsection{Gałąź pomocnicza Transformer}
Model z transformatorem wykorzystuje złożoną funkcję straty, łącząc ArcMargin i klasyfikację gałęzi pomocniczej.

\begin{table}[h]
\centering
\small
\begin{tabular}{ll}
\toprule
\textbf{Hiperparametr} & \textbf{Wartość} \\
\midrule
Rozmiar partii & 32 \\
Epoki & 25 \\
Optymalizator & SGD (momentum=0.9, weight\_decay=$5\times 10^{-4}$) \\
LR & $1 \times 10^{-2}$ \\
ALPHA (balans strat) & 0.4 \\
Transformer layers/heads & 6 / 8 \\
Wczesne zatrzymanie & Tak (cierpliwość=5) \\
Walidacja & 500 par z testu (średnia cos) \\
Augmentacja okluzji & $p=0.7$, wysokość 20 px \\
\bottomrule
\end{tabular}
\caption{Hiperparametry modelu z gałęzią Transformer.}
\label{tab:transformer_config}
\end{table}

\subsubsection{CBAM\_Block (uwaga w każdym bloku)}
Trening realizowany jest dwuetapowo z zamrożeniem backbone'u w pierwszej epoce.

\begin{table}[h]
\centering
\small
\begin{tabular}{ll}
\toprule
\textbf{Hiperparametr} & \textbf{Wartość} \\
\midrule
Rozmiar partii & 64 \\
Epoki & 30 \\
Optymalizator & SGD (momentum=0.9, weight\_decay=$5\times 10^{-4}$) \\
LR start & 0.1 \\
Harmonogram LR & StepLR (step\_size=10, gamma=0.1) \\
Zamrażanie & Epoka 1: tylko CBAM + head \\
Walidacja & 500 par z testu (średnia cos) \\
Augmentacja okluzji & $p=0.5$, wysokość 20 px \\
\bottomrule
\end{tabular}
\caption{Hiperparametry modelu CBAM\_Block.}
\label{tab:cbam_block_config}
\end{table}